\section{Technické nároky na inference}
\label{sec:selection-technical}

Tato sekce se věnuje technickým aspektům spouštění jazykových modelů, konkrétně rozdílu mezi využitím GPU a CPU, možnostem paralelizace a principu dávkového zpracování (batchování).
Pochopení těchto principů je důležité pro správné dimenzování on-premise infrastruktury a pro odhad propustnosti systému při různé zátěži.

\subsection{Rozdíl mezi použitím GPU a CPU}
\label{subsec:technical-gpu-vs-cpu}

Inference jazykového modelu je výpočetně náročná operace, jejíž jádro tvoří opakované násobení matic velkých rozměrů.
Pro každý generovaný token musí model provést tento výpočet skrze desítky nebo stovky vrstev architektury Transformer.
Způsob, na jakém hardware tento výpočet probíhá, zásadně ovlivňuje latenci (dobu zpracování jednoho dotazu) i celkovou propustnost systému.

\subsubsection*{Proč je GPU vhodný pro inferenci}

GPU (grafický procesor) jsou navrženy pro masivně paralelní výpočty.
Moderní spotřebitelský GPU jako NVIDIA RTX 4090 obsahuje 16\,384 specializovaných výpočetních jader (CUDA\footnote{CUDA (Compute Unified Device Architecture) je paralelní výpočetní platforma a programovací model vyvinutý firmou NVIDIA. Umožňuje využití GPU pro obecné výpočty.} jader) optimalizovaných pro operace s maticemi a vektory.
Výrobci GPU navíc implementují specializované hardwarové jednotky (Tensor Cores u NVIDIA), které dále akcelerují maticové výpočty potřebné při inferenci.

Klíčovým faktorem pro rychlost inference je šířka pásma paměti\footnote{Šířka pásma paměti (memory bandwidth) udává, kolik dat za sekundu může procesor přenést z paměti do výpočetních jader. Při inferenci je hlavním úzkým hrdlem právě tento přenos, nikoli samotná výpočetní rychlost.}.
GPU disponují typicky 500 až 2\,000 GB/s šířky pásma paměti, zatímco u serverových CPU je to 50 až 200 GB/s.
Protože největším úzkým hrdlem inference je právě přenos vah modelu z paměti do výpočetních jader, tato asymetrie přímo ovlivňuje rychlost generování.

Výsledek v praxi: model s 7 miliardami parametrů (7B) dokáže na moderním GPU typicky vygenerovat 20--80 tokenů za sekundu, zatímco na výkonném serverovém CPU je to zpravidla 1--5 tokenů za sekundu.

\subsubsection*{CPU inference a její omezení}

Provoz modelu na CPU je technicky možný, avšak výrazně pomalejší.
Nástroje jako \texttt{llama.cpp} \cite{llamacpp} jsou optimalizovány pro CPU inference a využívají SIMD\footnote{SIMD (Single Instruction, Multiple Data) je technika paralelizace na úrovni instrukční sady procesoru, která umožňuje provést jedinou instrukci na více datech zároveň. Moderní CPU podporují instrukční sady AVX-512, které rozšiřují šířku SIMD registrů na 512 bitů.}
instrukce pro urychlení maticových výpočtů.
Přesto lze tímto způsobem dosáhnout přijatelné rychlosti jen u velmi malých nebo silně kvantizovaných modelů.

Pro systém Kelvin, kde je analýza asynchronní a studenti nečekají na výsledek okamžitě, je CPU inference potenciálně přijatelná variantou při nízkém objemu odevzdání.
Při hromadném odevzdávání úloh (desítky odevzdání za hodinu) by fronty zpracování narůstaly na nepřijatelné hodnoty.

\begin{table}[H]
    \centering
    \resizebox{\textwidth}{!}{%
        \begin{tabular}{lll}
            \toprule
            \textbf{Parametr}           & \textbf{GPU (RTX 4090)}           & \textbf{CPU (16jádrový server)} \\
            \midrule
            Rychlost generování         & 30--80 tokenů/s (7B model)        & 2--6 tokenů/s (7B, Q4) \\
            Latence na 1 odevzdání      & 15--60 sekund                     & 3--15 minut \\
            Cena hardware               & $\sim$50\,000--80\,000 Kč          & Součást standardního serveru \\
            Paměť pro 7B model (Q4)     & 4--5 GB VRAM                      & 5--8 GB RAM \\
            Souběžné požadavky          & Batchování (viz~\ref{subsec:technical-batching}) & Prakticky bez batchování \\
            \bottomrule
        \end{tabular}
    }
    \caption{Orientační srovnání GPU a CPU inference pro model třídy 7B parametrů s kvantizací Q4}
    \label{tab:gpu-vs-cpu}
\end{table}

\subsection{Paralelizace zpracování}
\label{subsec:technical-parallelism}

V prostředí školního systému nepřicházejí požadavky na analýzu rovnoměrně.
Typickým scénářem je hromadné odevzdání těsně před termínem (deadlinem), kdy desítky studentů odešlou řešení v řádu minut.
V takovém případě vzniká fronta požadavků, které musí systém zpracovat.

Existují dvě základní strategie paralelizace zpracování:

\begin{enumerate}
    \item \textbf{Sekvenční zpracování s frontou} -- požadavky jsou zpracovávány postupně jeden po druhém.
    Tato varianta je nejjednodušší na implementaci, avšak studenti a učitelé musejí čekat, dokud nedojde na jejich požadavek.
    Pro nízké objemy odevzdání je tento přístup dostačující.

    \item \textbf{Paralelní zpracování} -- více požadavků je zpracováváno současně.
    U cloudového API to znamená odesílání více souběžných HTTP požadavků -- API poskytovatele zpracovává požadavky paralelně na svých serverech.
    U on-premise nasazení závisí na schopnostech inference serveru: \texttt{vLLM} \cite{vllm2023} je navržen tak, aby efektivně obsloužil desítky souběžných požadavků na jednom GPU.
\end{enumerate}

Pro systém Kelvin je klíčová skutečnost, že analýza je \textbf{asynchronní}: výsledky nejsou zobrazeny studentovi okamžitě, ale jsou k dispozici s určitým zpožděním (typicky v řádu minut).
To výrazně snižuje tlak na okamžitou propustnost a umožňuje použití jednoduchého front-based přístupu s nižší složitostí implementace.

\subsection{Batchování požadavků}
\label{subsec:technical-batching}

Batchování (dávkové zpracování) je technika, při níž jsou více požadavků zpracovány současně v jednom průchodu GPU.
Moderní GPU jsou schopny efektivně zpracovávat dávky matic najednou, díky čemuž se celkový čas potřebný pro zpracování více požadavků může být výrazně nižší než součet časů jednotlivých zpracování.

\subsubsection*{Jak batchování funguje}

Klíčová operace při inferenci -- maticové násobení pro attention výpočet -- je přirozeně paralelizovatelná přes dávku požadavků.
Místo toho, aby GPU provedl tuto operaci jednou pro jeden požadavek, provede ji jednou pro celou dávku, čímž efektivněji využije svou paralelní výpočetní kapacitu.
Výsledkem je vyšší propustnost (více tokenů za sekundu celkově) při mírném nárůstu latence jednotlivých požadavků.

\subsubsection*{Dynamické batchování ve vLLM}

Nástroj \texttt{vLLM} implementuje tzv. \emph{continuous batching} \cite{vllm2023} -- mechanismus, který dynamicky sestavuje dávky z příchozích požadavků bez nutnosti čekat na naplnění dávky pevné velikosti.
Díky tomu lze při hromadném odevzdávání výrazně zvýšit efektivitu využití GPU a zkrátit průměrnou dobu čekání.

\subsubsection*{Relevance pro Kelvin}

Pro systém Kelvin s asynchronním zpracováním jsou přínosy batchování nejvýraznější právě při hromadném odevzdávání.
Například pokud třicet studentů odevzdá řešení v průběhu pěti minut, server s vLLM může tyto požadavky zpracovat v dávkách po 4--8 souběžných analýzách, čímž se celková doba zpracování celé skupiny zkrátí přibližně třikrát až čtyřikrát oproti sekvenčnímu přístupu.

Při využití cloudového API je batchování interní záležitostí poskytovatele.
Provozovatel Kelvinu vidí pouze API endpoint -- o interním dávkování se nestará a efektivitu využití serveru řídí poskytovatel.

\endinput
